\section{Purpose}

\textbf{[STIPULATION]} This document treats the Casimir effect within the FFGFT
framework. The Casimir effect is in the legacy corpus the only experimentally
accessible bridge between the fundamental length scale $L_0=\xi\cdot\ell_P$ and
measurable vacuum forces. The document shows: (1) The modified Casimir formula
reproduces the standard formula exactly -- no new physics at currently accessible
separations. (2) The prefactor $4/3$ in $\xi=4/3\times10^{-4}$ is confirmed by
the Casimir effect independently of mass fits (A015). (3) There is a
characteristic vacuum length scale $L_\xi\approx100\,\mu\text{m}$ at which
Casimir energy density and a specific vacuum quantity coincide -- this is here
explicitly excluded.

\section{Fundamental Length Scales}

\textbf{[B]} FFGFT defines a hierarchy of characteristic length scales:
\[
  L_0 = \xi\cdot\ell_P \approx 2{.}155\times10^{-39}\,\text{m}
  \quad(\text{sub-Planck, minimal spacetime granulation}),
\]
\[
  \ell_P = \sqrt{\hbar G/c^3} \approx 1{.}616\times10^{-35}\,\text{m}
  \quad(\text{Planck length}),
\]
\[
  L_\xi \approx 100\,\mu\text{m}
  \quad(\text{characteristic vacuum length scale, measurable}).
\]
The coupling constant $\xi=4/3\times10^{-4}$ is the same parameter as in A020.
It can be derived from a Lagrangian describing the dynamics of the time field (A142).

\textbf{[B]} \textbf{Spacetime granulation.} At $L_0$ all vacuum fluctuations are
fully active. For $d>L_0$ only parts of these forces are visible -- this is the
physical origin of the $1/d^4$ dependence of the Casimir force. Due to the
extremely small size of $L_0\approx2\times10^{-39}\,\text{m}$, direct measurement
is not possible.

\section{The Modified Casimir Formula and Its Consistency}

\textbf{[K]} The modified Casimir energy density in the FFGFT framework:
\[
  |\rho_\text{Casimir}(d)|
  = \frac{\pi^2}{240\,\xi}\,\rho_\text{vac}\cdot\Bigl(\frac{L_\xi}{d}\Bigr)^4,
\]
where $\rho_\text{vac}=\xi\hbar c/L_\xi^4$ is a characteristic vacuum density.
Substituting:
\[
  |\rho_\text{Casimir}(d)|
  = \frac{\pi^2}{240\,\xi}\cdot\frac{\xi\hbar c}{L_\xi^4}\cdot\frac{L_\xi^4}{d^4}
  = \frac{\pi^2\hbar c}{240\,d^4}.
\]
This is exactly the standard Casimir formula. No contradiction, no new parameter --
the FFGFT formulation is numerically equivalent to the established QED results.

\textbf{[K]} \textbf{Numerical verification} (Casimir energy density vs.\ plate separation):

\begin{center}
\renewcommand{\arraystretch}{1.3}
{\small
\begin{tabular}{lcc}
\toprule
Separation $d$ & $\rho_\text{Casimir}$ (J/m³) & Ratio to $\rho_\text{vac}$ \\
\midrule
$100\,\mu$m & $4{.}17\times10^{-14}$ & 1{.}00 \\
$10\,\mu$m  & $4{.}17\times10^{-10}$ & $10^4$ \\
$1\,\mu$m   & $4{.}17\times10^{-2}$  & $10^{12}$ \\
\bottomrule
\end{tabular}
}
\renewcommand{\arraystretch}{1.0}
\end{center}

\textbf{[K]} \textbf{Characteristic scale hierarchy:}
\[
  \frac{L_0}{\ell_P} = \xi = 1{.}333\times10^{-4},
  \qquad
  \frac{\ell_P}{L_\xi} \approx 1{.}6\times10^{-31}.
\]

\section{Confirmation of the Factor $4/3$}

\textbf{[B]} The prefactor $4/3$ in $\xi=4/3\times10^{-4}$ appears independently
in the Casimir context: the factor $4/3$ appears in the Casimir energy term as
$E_\text{Casimir}\propto\tfrac43$ (the fourth -- musical interval $4{:}3$) and is
therefore not only a mass-fit quantity. This independence is recorded in A015 as
Casimir confirmation: no other prefactor (fifth $3/2$, third $5/4$) gives
consistent results for the mass spectrum \emph{and} the Casimir structure
simultaneously.

\section{Experimental Tests}

\textbf{[K]} Three classes of tests:

\textbf{Precision measurements at $d\approx L_\xi$:} At plate separation
$\approx100\,\mu$m the Casimir energy density reaches the order of $\rho_\text{vac}$.
This regime is experimentally accessible.

\textbf{Scaling behaviour:} The $1/d^4$ dependence should be exactly satisfied
down into the nanometre range. Deviations at $d\lesssim10\,\text{nm}$ could
indicate spacetime granulation.

\textbf{Experimental $L_\xi$ values from the literature:} Parallel plates:
$\sim228\,\text{nm}$; sphere-plate geometry: $\sim1{.}75\,\mu\text{m}$ to
$\sim18\,\mu\text{m}$. The scatter reflects geometric differences
($F\propto1/d^4$ vs.\ $F\propto1/d^3$).

\section{Verification Script}

\begin{itemize}[leftmargin=2em,itemsep=2pt,topsep=2pt]
  \item Numerical verification of the modified Casimir formula (reproduction of
    the standard formula), length scale hierarchy, comparison with $L_\xi$ values:
    \texttt{python/A\_Serie\_Skripte/a260\_casimir.py}
\end{itemize}

\section{Open Questions}

\textbf{Direct measurement of $L_0$ is not possible.} The minimal length scale
$L_0\approx2\times10^{-39}\,\text{m}$ lies far beyond any experimental reach.
The characteristic scale $L_\xi$ is experimentally accessible, but its precise
connection to $L_0$ presupposes the FFGFT framework.

\textbf{Deviations from the $1/d^4$ law at $d\sim10\,\text{nm}$ are not measured.}
The prediction of a scale change near the sub-Planck boundary is qualitative,
not quantitatively worked out.

\section{Supersedes}

This document subsumes: Dok.~091 (modified Casimir theory, length scale hierarchy,
numerical verification, experimental predictions and tests). \emph{Not} taken over:
the Casimir vacuum scaling from Dok.~091 (excluded sector, outside A-series core)
and the macroscopic vacuum implications.

\section{Use of AI Tools}

This document was created with the assistance of AI tools. Responsibility
for content and errors rests with the author. Full wording of the rules in A010.
